% Options for packages loaded elsewhere
% Options for packages loaded elsewhere
\PassOptionsToPackage{unicode}{hyperref}
\PassOptionsToPackage{hyphens}{url}
%
\documentclass[
  11pt,
  letterpaper,
  DIV=11,
  numbers=noendperiod,
  twoside]{scrartcl}
\usepackage{xcolor}
\usepackage[left=1.1in, right=1in, top=0.8in, bottom=0.8in,
paperheight=9.5in, paperwidth=7in, includemp=TRUE, marginparwidth=0in,
marginparsep=0in]{geometry}
\usepackage{amsmath,amssymb}
\setcounter{secnumdepth}{3}
\usepackage{iftex}
\ifPDFTeX
  \usepackage[T1]{fontenc}
  \usepackage[utf8]{inputenc}
  \usepackage{textcomp} % provide euro and other symbols
\else % if luatex or xetex
  \usepackage{unicode-math} % this also loads fontspec
  \defaultfontfeatures{Scale=MatchLowercase}
  \defaultfontfeatures[\rmfamily]{Ligatures=TeX,Scale=1}
\fi
\usepackage{lmodern}
\ifPDFTeX\else
  % xetex/luatex font selection
  \setmathfont[]{Garamond-Math}
\fi
% Use upquote if available, for straight quotes in verbatim environments
\IfFileExists{upquote.sty}{\usepackage{upquote}}{}
\IfFileExists{microtype.sty}{% use microtype if available
  \usepackage[]{microtype}
  \UseMicrotypeSet[protrusion]{basicmath} % disable protrusion for tt fonts
}{}
\usepackage{setspace}
% Make \paragraph and \subparagraph free-standing
\makeatletter
\ifx\paragraph\undefined\else
  \let\oldparagraph\paragraph
  \renewcommand{\paragraph}{
    \@ifstar
      \xxxParagraphStar
      \xxxParagraphNoStar
  }
  \newcommand{\xxxParagraphStar}[1]{\oldparagraph*{#1}\mbox{}}
  \newcommand{\xxxParagraphNoStar}[1]{\oldparagraph{#1}\mbox{}}
\fi
\ifx\subparagraph\undefined\else
  \let\oldsubparagraph\subparagraph
  \renewcommand{\subparagraph}{
    \@ifstar
      \xxxSubParagraphStar
      \xxxSubParagraphNoStar
  }
  \newcommand{\xxxSubParagraphStar}[1]{\oldsubparagraph*{#1}\mbox{}}
  \newcommand{\xxxSubParagraphNoStar}[1]{\oldsubparagraph{#1}\mbox{}}
\fi
\makeatother


\usepackage{longtable,booktabs,array}
\newcounter{none} % for unnumbered tables
\usepackage{calc} % for calculating minipage widths
% Correct order of tables after \paragraph or \subparagraph
\usepackage{etoolbox}
\makeatletter
\patchcmd\longtable{\par}{\if@noskipsec\mbox{}\fi\par}{}{}
\makeatother
% Allow footnotes in longtable head/foot
\IfFileExists{footnotehyper.sty}{\usepackage{footnotehyper}}{\usepackage{footnote}}
\makesavenoteenv{longtable}
\usepackage{graphicx}
\makeatletter
\newsavebox\pandoc@box
\newcommand*\pandocbounded[1]{% scales image to fit in text height/width
  \sbox\pandoc@box{#1}%
  \Gscale@div\@tempa{\textheight}{\dimexpr\ht\pandoc@box+\dp\pandoc@box\relax}%
  \Gscale@div\@tempb{\linewidth}{\wd\pandoc@box}%
  \ifdim\@tempb\p@<\@tempa\p@\let\@tempa\@tempb\fi% select the smaller of both
  \ifdim\@tempa\p@<\p@\scalebox{\@tempa}{\usebox\pandoc@box}%
  \else\usebox{\pandoc@box}%
  \fi%
}
% Set default figure placement to htbp
\def\fps@figure{htbp}
\makeatother


% definitions for citeproc citations
\NewDocumentCommand\citeproctext{}{}
\NewDocumentCommand\citeproc{mm}{%
  \begingroup\def\citeproctext{#2}\cite{#1}\endgroup}
\makeatletter
 % allow citations to break across lines
 \let\@cite@ofmt\@firstofone
 % avoid brackets around text for \cite:
 \def\@biblabel#1{}
 \def\@cite#1#2{{#1\if@tempswa , #2\fi}}
\makeatother
\newlength{\cslhangindent}
\setlength{\cslhangindent}{1.5em}
\newlength{\csllabelwidth}
\setlength{\csllabelwidth}{3em}
\newenvironment{CSLReferences}[2] % #1 hanging-indent, #2 entry-spacing
 {\begin{list}{}{%
  \setlength{\itemindent}{0pt}
  \setlength{\leftmargin}{0pt}
  \setlength{\parsep}{0pt}
  % turn on hanging indent if param 1 is 1
  \ifodd #1
   \setlength{\leftmargin}{\cslhangindent}
   \setlength{\itemindent}{-1\cslhangindent}
  \fi
  % set entry spacing
  \setlength{\itemsep}{#2\baselineskip}}}
 {\end{list}}
\usepackage{calc}
\newcommand{\CSLBlock}[1]{\hfill\break\parbox[t]{\linewidth}{\strut\ignorespaces#1\strut}}
\newcommand{\CSLLeftMargin}[1]{\parbox[t]{\csllabelwidth}{\strut#1\strut}}
\newcommand{\CSLRightInline}[1]{\parbox[t]{\linewidth - \csllabelwidth}{\strut#1\strut}}
\newcommand{\CSLIndent}[1]{\hspace{\cslhangindent}#1}



\setlength{\emergencystretch}{3em} % prevent overfull lines

\providecommand{\tightlist}{%
  \setlength{\itemsep}{0pt}\setlength{\parskip}{0pt}}



 


\setlength\heavyrulewidth{0ex}
\setlength\lightrulewidth{0ex}
\usepackage[automark]{scrlayer-scrpage}
\clearpairofpagestyles
\cehead{
  Brian Weatherson
  }
\cohead{
  Fairness Has Memory
  }
\ohead{\bfseries \pagemark}
\cfoot{}
\makeatletter
\newcommand*\NoIndentAfterEnv[1]{%
  \AfterEndEnvironment{#1}{\par\@afterindentfalse\@afterheading}}
\makeatother
\NoIndentAfterEnv{itemize}
\NoIndentAfterEnv{enumerate}
\NoIndentAfterEnv{description}
\NoIndentAfterEnv{quote}
\NoIndentAfterEnv{equation}
\NoIndentAfterEnv{longtable}
\NoIndentAfterEnv{abstract}
\renewenvironment{abstract}
 {\vspace{-1.25cm}
 \quotation\small\noindent\emph{Abstract}:}
 {\endquotation}
\newfontfamily\tfont{EB Garamond}
\addtokomafont{disposition}{\rmfamily}
\addtokomafont{descriptionlabel}{\rmfamily}
\addtokomafont{title}{\normalfont\itshape}
\let\footnoterule\relax

\makeatletter
\renewcommand{\@maketitle}{%
  \newpage
  \null
  \vskip 2em%
  \begin{center}%
  \let \footnote \thanks
    {\itshape\huge\@title \par}%
    \vskip 0.5em%  % Reduced from default
    {\large
      \lineskip 0.3em%  % Reduced from default 0.5em
      \begin{tabular}[t]{c}%
        \@author
      \end{tabular}\par}%
    \vskip 0.5em%  % Reduced from default
    {\large \@date}%
  \end{center}%
  \par
  }
\makeatother
\RequirePackage{lettrine}

\renewenvironment{abstract}
 {\quotation\small\noindent\emph{Abstract}:}
 {\endquotation\vspace{-0.02cm}}

\setmainfont{EB Garamond Math}[
  BoldFont = {EB Garamond SemiBold},
  ItalicFont = {EB Garamond Italic},
  RawFeature = {+smcp},
]

\newfontfamily\scfont{EB Garamond Regular}[RawFeature=+smcp]
\renewcommand{\textsc}[1]{{\scfont #1}}

\renewcommand{\LettrineTextFont}{\scfont}
% Number theorem-like environments straight through (Theorem 1, Theorem 2, ...)
% rather than by section (1.1, 1.2, ...), which is Quarto's LaTeX default.
% Guarded, since Quarto only defines a counter for environments the document uses.
\makeatletter
\AtBeginDocument{%
  \@for\bw@thm:={theorem,lemma,corollary,proposition,conjecture,definition,%
                 example,exercise,refremark,refsolution}\do{%
    \@ifundefined{c@\bw@thm}{}{\expandafter\counterwithout\expandafter{\bw@thm}{section}}%
  }%
}
\makeatother
\KOMAoption{captions}{tableheading}
\makeatletter
\@ifpackageloaded{caption}{}{\usepackage{caption}}
\AtBeginDocument{%
\ifdefined\contentsname
  \renewcommand*\contentsname{Table of contents}
\else
  \newcommand\contentsname{Table of contents}
\fi
\ifdefined\listfigurename
  \renewcommand*\listfigurename{List of Figures}
\else
  \newcommand\listfigurename{List of Figures}
\fi
\ifdefined\listtablename
  \renewcommand*\listtablename{List of Tables}
\else
  \newcommand\listtablename{List of Tables}
\fi
\ifdefined\figurename
  \renewcommand*\figurename{Figure}
\else
  \newcommand\figurename{Figure}
\fi
\ifdefined\tablename
  \renewcommand*\tablename{Table}
\else
  \newcommand\tablename{Table}
\fi
}
\@ifpackageloaded{float}{}{\usepackage{float}}
\floatstyle{ruled}
\@ifundefined{c@chapter}{\newfloat{codelisting}{h}{lop}}{\newfloat{codelisting}{h}{lop}[chapter]}
\floatname{codelisting}{Listing}
\newcommand*\listoflistings{\listof{codelisting}{List of Listings}}
\makeatother
\makeatletter
\makeatother
\makeatletter
\@ifpackageloaded{caption}{}{\usepackage{caption}}
\@ifpackageloaded{subcaption}{}{\usepackage{subcaption}}
\makeatother
\usepackage{bookmark}
\IfFileExists{xurl.sty}{\usepackage{xurl}}{} % add URL line breaks if available
\urlstyle{same}
% fallback for those not using the hyperref driver hyperxmp:
\makeatletter
\@ifundefined{xmpquote}{\newcommand{\xmpquote}[1]{#1}}{}
\makeatother
\hypersetup{
  pdftitle={Fairness Has Memory},
  pdfauthor={Claude},
  hidelinks,
  pdfcreator={LaTeX via pandoc}}


\title{Fairness Has Memory}
\author{Claude}
\date{2026}
\begin{document}
\maketitle
\begin{abstract}
When a scarce, indivisible good is allocated by lottery among people
with equal claims, the standard theory of fairness says the lottery is
the fairest thing available. When the same kind of good is allocated the
same way year after year, to largely the same people, the standard
theory has nothing further to say: each year's lottery is fair, so the
sequence is fair. I argue that this is a mistake. A sequence of
individually fair lotteries can be an unfair allocation, and fairness
requires that a claimant's chances in later rounds depend on the history
of her unsatisfied claims. The argument uses only premises the standard
theory already accepts: that a divisible good must be divided in
proportion to claims, and that a lottery is a second-best surrogate for
division. Over a sequence, the good becomes divisible. I then show what
memory should track and how heavily it may press. It tracks a claimant's
shortfall from her fair share, in units of the good where later units
are as good as earlier ones and in waiting where they are not; and it
may press that shortfall only as hard as is consistent with every
claimant's still receiving, in each period, the surrogate she is owed.
Rules that forget, like New Mexico's hunting draw and the United States
visa lotteries, and rules that remember lexically, like preference-point
queues, are both unfair, in opposite directions. Finally I say how far
back, and how far across, fair memory reaches.
\end{abstract}


\setstretch{1.1}
\section{Two lotteries}\label{sec-intro}

Every year Nevada and New Mexico each have more people who want to hunt
bighorn sheep than they have sheep to spare, so each holds a lottery.
The two lotteries differ in one respect that will occupy this paper.
Nevada remembers. An applicant who is unsuccessful earns a bonus point;
her points are squared to fix how many random numbers are drawn for her
the next year, and the lowest number she draws is her entry, so ten
years of losing gives her a hundred and one draws against a newcomer's
one.\footnote{Nevada Department of Wildlife, ``Bonus Point Program,''
  https://www.ndow.org/blog/bonus-point-program/. Points revert to zero
  when the applicant draws a tag, or after two consecutive years of not
  applying.} New Mexico forgets. Its Department of Game and Fish states
the policy in one sentence: ``New Mexico does NOT grant preference to
applicants who were unsuccessful in previous drawings.''\footnote{New
  Mexico Department of Game and Fish, ``How the New Mexico Draw Works,''
  https://wildlife.dgf.nm.gov/hunting/applications-and-draw-information/how-new-mexico-draw-works/.}
A hunter passed over for ten years enters the eleventh draw with exactly
the chances of someone entering for the first time.

Which of these is fair? Most people, asked, side with Nevada. When
Huesch and Brady (\citeproc{ref-HueschBrady2010}{2010}) asked subjects
to design a pair of back-to-back raffles among the same ten people, 78
per cent said the winner of the first should be excluded from the
second, and 75 per cent said so even when asked as disinterested
spectators. Many institutions side with Nevada too, though they disagree
about how to remember. Arizona gives an unsuccessful hunter one extra
random number per year rather than squaring them. In Colorado's and
Wyoming's main draws, licences go to whoever has accumulated the most
preference points, and chance only breaks ties. The Grand Canyon's
river-permit lottery gives an applicant one ticket for each year since
she last held a permit, to a maximum of five. Kidney allocation in the
United States is a points system whose principal component is waiting
time. Publicly funded IVF in England is capped at a fixed number of
cycles per person.\footnote{Colorado Parks and Wildlife, ``Primary
  Draw,'' https://cpw.state.co.us/hunting/big-game/primary-draw; Wyoming
  Game and Fish Department, ``Understanding the Draw'' (2024); Arizona
  Game and Fish Department, ``Bonus Point Process''; Grand Canyon
  National Park, ``Noncommercial River Trip FAQ,''
  https://grcariverpermits.nps.gov/faq.cfm; National Academies of
  Sciences, Engineering, and Medicine, \emph{Realizing the Promise of
  Equity in the Organ Transplantation System} (2022), ch.~5; NICE
  guideline NG257, ``Fertility problems: assessment and treatment''
  (2026), recommendations 1.39.6--1.39.9. The NBA draft lottery,
  weighted by the previous season's record, and the Amsterdam housing
  register, which awards a point per year registered and a further point
  per month of unsuccessful applications, also remember; the Wimbledon
  public ballot and most American school-admission lotteries do not.}
And then there are the systems that forget: New Mexico, the Wimbledon
ballot, and, most consequentially, the lotteries by which the United
States allocates the roughly 85,000 H-1B visas Congress permits each
year and the 55,000 Diversity Visas. A skilled worker whose registration
was not selected in three successive years enters the fourth draw with
the chances of a newcomer. That remains true under the rule that took
effect in February 2026, which weights H-1B registrations by the wage
level of the job offered, entering the highest-paid four times and the
lowest once, but gives no weight to how many times a person has already
been passed over.\footnote{Department of Homeland Security, ``Weighted
  Selection Process for Registrants and Petitioners Seeking To File
  Cap-Subject H-1B Petitions,'' 90 \emph{Federal Register} 60864 (29
  December 2025), effective 27 February 2026.}

The theory of fair allocation has not connected this disagreement to the
theory of claims, and where it touches it, it sides with New Mexico. The
best-developed account, due to John Broome
(\citeproc{ref-Broome1984}{1984a}, \citeproc{ref-Broome1990}{1991a},
\citeproc{ref-Broome1991}{1991b}), holds that fairness requires people's
claims to a good to be satisfied in proportion to their strength, and
that when the good is indivisible and claims are equal, an equal lottery
is the fairest available procedure. The theory evaluates a lottery on
its own. On its own, each year's New Mexico draw is impeccable, and if
the sequence is nothing but its members, the sequence is impeccable too.
Broome's critics, who deny that chances are a kind of satisfaction and
locate the fairness of lotteries in what they express or exclude
(\citeproc{ref-Wasserman1996}{Wasserman 1996};
\citeproc{ref-Stone2007}{Stone 2007}, \citeproc{ref-Stone2011}{2011};
\citeproc{ref-Spiekermann2022}{Spiekermann 2022}), reach the same
verdict by a different route: an impartial procedure that treats
everyone identically is fair, whatever last year's procedure did. The
wider literature on lotteries (\citeproc{ref-Sher1980}{Sher 1980};
\citeproc{ref-KornhauserSager1988}{Kornhauser and Sager 1988};
\citeproc{ref-Elster1989}{Elster 1989};
\citeproc{ref-Duxbury1999}{Duxbury 1999};
\citeproc{ref-Saunders2008}{Saunders 2008};
\citeproc{ref-Stone2008}{Stone 2008}) takes the single lottery as its
unit; Goodwin's (\citeproc{ref-Goodwin2005}{2005}) ``total social
lottery,'' which periodically redistributes society's goods, is the
forgetting rule adopted as a constitution. The recent literature on
queues (\citeproc{ref-JohnMillum2020}{John and Millum 2020};
\citeproc{ref-HerschRowe2024}{Hersch and Rowe 2024}) asks whether
waiting should count, which is a question about history, but does not
connect the answer to claims. The nearest things to the thesis I will
defend are two passages by Hersch and Rowe. In one, a parent hands out a
chocolate bar every ten minutes, and the child who got the first insists
that ``chocolate bar at t1 is not the same good as chocolate bar at t2,
and a lottery must be had for this new and distinct good''; they reply
that if we index goods to times ``we lose something that matters for
fairness,'' without saying what (\citeproc{ref-HerschRowe2024}{Hersch
and Rowe 2024, sec. 5}). In the other, a survey, they observe that
recurring chores are ``not independent of each other,'' that a recurring
chore ``which extends temporally'' is not indivisible even when each
instance is, and that ``who completed the chore last time matters to who
has a claim to avoid it this time''
(\citeproc{ref-HerschRowe2025}{Hersch and Rowe 2025, 3}). Both
observations are right. This paper says what is lost, why, how far it
reaches, and what follows for the institutions just described.

The thesis is that fairness has memory. More precisely, I will argue for
three claims.

\emph{Non-separability.} The fairness of a sequence of allocations is
not a function of the fairness of its members as the standard theory
assesses them. A sequence of lotteries each of which is fair, on that
theory, can be and typically is an unfair allocation.

\emph{Memory.} When a good is allocated repeatedly among overlapping
claimants, fairness requires that each claimant's chances in later
rounds depend on her history of unsatisfied claims.

\emph{Currency and weight.} What memory should track is a claimant's
shortfall from her fair share, in the currency in which her claim is
divisible: units of the good, where a later unit is as good to her as an
earlier one; waiting, where it is not. And it may press that shortfall
only as hard as is consistent with every claimant's still receiving, in
each period, the surrogate she is owed. A rule that forgets fails the
first requirement. A rule that lets history decide, so that a newcomer's
present claim counts for nothing, fails the second. Both are in use.

The argument for the first two claims (Section~\ref{sec-division}) needs
nothing beyond premises that Broome and most of his critics already
accept: that when a good can be divided among equal claimants, fairness
requires that it be divided, and that a lottery for an indivisible good
is an admittedly imperfect surrogate for division. Over a sequence, the
good becomes divisible. The argument seems to conflict with the standard
theory's doctrine that losing a fair lottery extinguishes one's claim;
Section~\ref{sec-extinction} shows that the conflict is only apparent.
Section~\ref{sec-proportion} turns to the third claim and to design.
Section~\ref{sec-reach} asks how far back, and how far across, memory
reaches: does a housing loss count in a visa lottery, a parent's loss
for a child, bad luck at the hands of no one at all? The answers are
derived rather than stipulated, and they keep the thesis clear of luck
egalitarianism and of theories of rectification.
Section~\ref{sec-objections} answers the remaining objections.

\section{Claims, chances, and division}\label{sec-broome}

Broome's theory begins with claims. A claim, in his sense, is a duty
owed to the candidate herself that she should have the good, as distinct
from other reasons for giving it to her, such as that doing so would be
useful (\citeproc{ref-Broome1990}{Broome 1991a, 90--92}). Need is the
paradigm source of claims; desert and prior agreement are others.
Fairness is a specific requirement on how claims are treated: ``equal
claims require equal satisfaction,'' ``stronger claims require more
satisfaction than weaker ones,'' and ``weaker claims require some
satisfaction. Weaker claims must not simply be overridden by stronger
ones'' (\citeproc{ref-Broome1990}{Broome 1991a, 95}). Fairness is thus
comparative, ``concerned only with relative satisfaction, not absolute
satisfaction'' (\citeproc{ref-Broome1990}{Broome 1991a, 95}); it can be
perfectly met by denying the good to everyone, and the requirement that
claims be satisfied as much as possible competes with it
(\citeproc{ref-Broome1990}{Broome 1991a, 97}).

Two features of the theory will bear the weight of my argument. The
first is what I will call the Division Principle. Where a good is
divisible, proportional satisfaction is straightforwardly achievable,
and the theory says it is required. If two people have equal claims to a
quantity of food, fairness says the food ``should be divided,'' even
when giving it all to one of them would do more good
(\citeproc{ref-Broome1990}{Broome 1991a, 95}). Broome states the
principle at its most general in a passage that already contains, in
miniature, the idea that a good can be divided across time:

\begin{quote}
I have been implicitly assuming that the good to be distributed is
indivisible. Sometimes we are faced with an indivisibility in the nature
of things; the good, for instance, of staying in a lifeboat cannot be
divided. But sometimes the problem of selection arises even with a good
that could be divided. The good, for instance, of avoiding military
service could be divided by conscripting everyone for a short time
rather than a few people for a long time. If people's claims to a good
are equal, then fairness requires that the good should be divided
between them if it can be. That is the only way to treat them with
perfect equality. (\citeproc{ref-Broome1984}{Broome 1984a, 48})
\end{quote}

The second feature is the account of lotteries as a second best. When
the good really is indivisible, and there is not enough of it, ``some
unfairness is inevitable,'' because ``some candidates will get the good
and others will not'' (\citeproc{ref-Broome1990}{Broome 1991a, 97--98}).
What a lottery achieves is ``a sort of partial equality in
satisfaction'': ``each person can be given a sort of surrogate
satisfaction. By holding a lottery, each can be given an equal chance of
getting the good. This is not a perfect fairness, but it meets the
requirement of fairness to some extent''
(\citeproc{ref-Broome1990}{Broome 1991a, 98}). A chance counts as
surrogate satisfaction because ``if you have a chance of getting the
good you may actually get it,'' unlike merely giving a claim its due
weight in a deliberation that goes against it
(\citeproc{ref-Broome1990}{Broome 1991a, 98}). And the surrogate is
explicitly a surrogate: the result of a lottery among unequal claimants
``will generally be that the good goes to candidates who do not have the
strongest claims. This is less fair than the result of giving it
directly to those who do'' (\citeproc{ref-Broome1990}{Broome 1991a,
98--99}). In the 1984 paper the lottery achieves ``a fairness that is
less than perfect,'' and is what one does when division, the perfectly
fair option, ``may be very damaging to the general good''
(\citeproc{ref-Broome1984}{Broome 1984a, 48}). Two things follow that I
will use. Where division is available, it is required. Where it is not,
the surrogate is not merely permitted but owed: a claimant who cannot be
given her portion must be given her chance.

The theory has been much criticized (\citeproc{ref-Hooker2005}{Hooker
2005}; \citeproc{ref-Tomlin2012}{Tomlin 2012};
\citeproc{ref-Lazenby2014}{Lazenby 2014};
\citeproc{ref-KirkpatrickEastwood2015}{Kirkpatrick and Eastwood 2015};
\citeproc{ref-Piller2017}{Piller 2017}), and weighted lotteries have
been refined and contested chiefly in the debate about whether the
numbers count (\citeproc{ref-Kamm1993}{Kamm 1993};
\citeproc{ref-Broome1998}{Broome 1998};
\citeproc{ref-Timmermann2004}{Timmermann 2004};
\citeproc{ref-Henning2015}{Henning 2015};
\citeproc{ref-Saunders2010}{Saunders 2010}; \citeproc{ref-Vong2020}{Vong
2020}; \citeproc{ref-Tank2024}{Tank, Wendler, and Carstensen-Mainka
2024}). But the Division Principle and the second-best status of
lotteries are common ground. Wasserman
(\citeproc{ref-Wasserman1996}{1996}) doubts that chances are a kind of
satisfaction, since ``hungry claimants'' cannot ``actually eat chances''
(p.~47), and holds that among equally entitled claimants ``we have no
reason save economy to equalize their expected rather than their actual
outcomes'' (p.~46); he also counts ``time sharing'' as literal division
of a good (p.~30 n.~1). Stone (\citeproc{ref-Stone2007}{2007, 284--85};
\citeproc{ref-Stone2011}{2011}) holds that a lottery is just because
impartiality demands that some be favoured over others ``only for
legitimate reasons,'' and a lottery keeps illegitimate ones out of a
decision no legitimate reason can settle; but dividing a divisible good
is not such a decision. Lazenby (\citeproc{ref-Lazenby2014}{2014, 341})
holds that whatever fairness a lottery contributes is lexically
posterior to fairness in outcome. My argument needs only what these
positions share.

One more preliminary. The theory speaks of claims, and it may be asked
what duty Nevada owes anyone to let her shoot a sheep. The answer is
that the claims in the cases that concern me are claims from the
allocator's own undertaking. When a public body announces that it will
allocate a public resource among eligible applicants by a fair
procedure, it undertakes to each applicant that her application will be
treated as the procedure promises; each eligible applicant thereby has a
claim, owed to her, that is equal to every other's unless the rules
distinguish them. Agreement is one of Broome's own sources of claims,
and this is agreement of a public kind. It extends to a private
allocator that advertises a ballot, and it is what makes the visa
registrant's position a matter of fairness rather than of policy alone.
What it does not do is turn every reason for allocating a visa into a
claim. The wage level of the job an H-1B beneficiary has been offered is
a reason of the other kind, a reason of usefulness, and the 2026 reform
that weights the lottery by it has made the lottery more responsive to
usefulness, not to claims. I return to this in
Section~\ref{sec-bottleneck}.

\section{The divisibility argument}\label{sec-division}

\subsection{The setting}\label{the-setting}

Consider a good that comes in indivisible units, one unit per period,
allocated by a single agency among a set of claimants whose claims are
equal, and where receiving a unit in one period does not remove a
claimant from the pool: a bighorn tag a year, a river launch a season, a
night off the rota a week. To begin with I assume the claimants are the
same from period to period; I relax this in
Section~\ref{sec-proportion}. Let there be \emph{N} claimants and
\emph{T} periods. Within any one period the standard theory's verdict is
clear and, I think, correct: one unit, \emph{N} equal claims, no way to
divide the unit, so hold an equal lottery. The question is what the
theory says about the \emph{T} periods taken together.

\subsection{The argument}\label{the-argument}

The units are indivisible. The sequence of units is not. Over \emph{T}
periods, it is possible to give each claimant \emph{T}/\emph{N} units,
or as near to that as the integers allow: if \emph{T} is a multiple of
\emph{N}, exactly \emph{T}/\emph{N} each; otherwise some get one more
than others, but nobody's total differs from anybody else's by more than
one. Call the procedure that achieves this \emph{rotation}: draw a
random order once, and then cycle through it.\footnote{When \emph{T} is
  not a multiple of \emph{N}, the random initial order settles fairly
  who gets the extra unit. Rotation is not the only procedure with this
  property; excluding past winners until everyone has won and then
  re-drawing has the same distribution of totals. Nothing below turns on
  the difference.} Rotation divides the sequence among the claimants in
proportion to their equal claims. It is, in the sequence, what dividing
the food is in the single case.

The Division Principle says that ``if people's claims to a good are
equal, then fairness requires that the good should be divided between
them if it can be.'' The sequence can be divided. So fairness requires
that it be divided. So fairness requires rotation, or something that
comes as close to it as circumstances allow.

Now compare the procedure New Mexico follows: a fresh, independent,
equal lottery in every period. Call it the \emph{forgetting} rule. In
any one period it does what the theory asks. Over \emph{T} periods it
does not divide the sequence; it raffles it. With ten claimants and ten
periods, the forgetting rule leaves, on average, three and a half of the
ten with nothing at all, while two or three of the others receive two or
more units. With a hundred claimants and a hundred periods, 37 of the
hundred will never have won and 26 will have won twice or
more.\footnote{The expected number of claimants with nothing is
  \(N(1-1/N)^T\); for \(N = T = 100\) this is
  \(100 \times 0.99^{100} \approx 36.6\). The other figures are computed
  in the same way, or by simulation as described in the Appendix.}
Rotation leaves nobody with nothing. The forgetting rule does over the
sequence exactly what the theory forbids in the single case: it takes a
divisible good and gives it out by lot.

Here is the point put as a reductio. Suppose someone proposed to
allocate a cake among ten people with equal claims as follows: cut it
into ten slices, and raffle each slice separately. Each raffle is, taken
by itself, exactly what the theory recommends for an indivisible good.
Yet nobody thinks this is a fair way to allocate a cake. The theory
itself says the cake should be divided, and dividing it means giving
everyone a slice, not giving everyone a chance at each slice. The only
difference between the cake and the tags is that the slices are
allocated simultaneously and the tags a year apart. If simultaneity is
not what makes the cake divisible, and it is not, since a cake sliced
today and handed out over ten days is still a divisible cake, then the
tags are divisible too.

Two consequences follow. The first is the failure of separability. The
standard theory assesses a lottery by the claims present when it is
held, and it takes those claims to be fixed by need, desert, agreement
and the like, not by what earlier lotteries did. Say that an allocation
rule is \emph{locally fair} at a period if the procedure it uses at that
period, assessed in that way, is fair. The forgetting rule is locally
fair at every period: at each draw, everyone present has an equal claim
and an equal chance. Rotation is not locally fair at any period after
the first, since given the history one claimant is certain to receive
the unit and the others certain not to. Yet rotation is the fair
allocation of the sequence and the forgetting rule is not. So local
fairness at every period is neither necessary nor sufficient for the
fairness of the sequence: fairness, assessed as the standard theory
assesses it, is not separable across occasions. There is another way to
put the same result, and from Section~\ref{sec-extinction} on I will use
it. If one lets claims carry history, so that a claimant who has
received less than her share has, on that account, a stronger claim to
the next unit, then rotation \emph{is} locally fair at every period and
the forgetting rule is not, and fairness is separable after all.
Non-separability of fairness with respect to history-blind claims and
memory in claims are one thesis in two vocabularies. What is not
available is the standard theory's combination: history-blind claims and
a separable notion of fairness.\footnote{Broome's theory of goodness is
  organized around separability principles for betterness across people,
  times, and states (\citeproc{ref-Broome1991}{Broome 1991b}).
  Non-separability of fairness across occasions is a different kind of
  non-additivity, of a comparative pattern rather than a value, and I do
  not claim it bears on those principles. It does bear on how fairness
  is weighed against goodness: if fairness attaches to sequences, the
  weighing must be done at the level of the sequence.}

The second consequence is memory. Rotation is a rule under which who
receives the unit in period \emph{t} depends on who received it in
periods before \emph{t}. Any rule that divides the sequence must have
this feature, since dividing the sequence means giving units to those
who have received fewer, and that requires knowing who has received
fewer. So if fairness requires division of the sequence, fairness
requires memory.

\subsection{What the argument assumes}\label{sec-assumes}

The argument does not assume that chances are a kind of satisfaction. It
assumes only that a lottery is a second best relative to division, which
is common ground.

It does assume something that needs stating, because it is where a
careful opponent will press. The Division Principle applies to a good
that ``could be divided.'' In Broome's examples the claimants end up
with portions of one thing, and the tags a hunter draws in different
years are not portions of one tag. What licenses treating them as
portions of one good is a premise about the claim. Call it
\emph{Fungibility}: units allocated in different periods are portions of
one divisible good, for a claimant, to the extent that her claim is to a
unit of that kind and a later unit would satisfy it as an earlier one
would. The hunter's claim is to a sheep tag, not to the 2027 tag; the
conscript's, in Broome's example, is to avoid service, not to avoid it
this month. Where Fungibility holds, the sequence is divisible in the
currency of the claims, not merely in the allocator's bookkeeping.

Fungibility is a matter of degree, because delay is nearly always some
cost, and it fails outright in two ways. It fails when the claim is to a
particular unit: the one transplant organ that matches, the seat on the
one lifeboat. Then there is no sequence and the standard theory applies
as it stands. And it fails, more interestingly, when the claim is to a
unit of the kind but a later unit is much worse for the claimant: the
housing applicant who will be evicted before next year's allocation, the
visa applicant whose sponsor will not wait. Here the claim is largely a
claim to timeliness. But timeliness is divisible: delay can be shared
out among claimants as units cannot. In this case the argument does not
lapse; it changes currency. Fairness requires that the \emph{waiting} be
divided in proportion to claims, and the memory this requires tracks
time waited (Section~\ref{sec-bottleneck}). In the general case, where a
later unit is worth something but less, the currency is units discounted
by delay, and the pure cases are its endpoints.

The argument does not assume equal claims either. With unequal claims
the sequence should be divided in proportion to them, and doing so with
indivisible units is the apportionment problem; the divisor methods that
solve it for parliaments have a sequential form, which is a rule with
memory.\footnote{Wintein and Heilmann
  (\citeproc{ref-WinteinHeilmann2018}{2018}) show that the divisor
  methods of apportionment theory divide several indivisible units
  proportionally without lotteries; the classic treatment is Balinski
  and Young (\citeproc{ref-BalinskiYoung2001}{2001}). In sequential form
  a divisor method gives each unit to the claimant who maximizes the
  ratio of her claim to the number of units she has already received
  plus a constant; with equal claims every divisor method reduces to
  rotation. The same point has lately been made in computer science,
  where ``repeated'' and ``temporal'' fair division begins from the
  observation that an allocation of indivisible items that cannot be
  made envy-free in one round can be made envy-free over enough rounds
  (\citeproc{ref-Igarashi2024}{Igarashi et al. 2024};
  \citeproc{ref-AleksandrovWalsh2020}{Aleksandrov and Walsh 2020}).}

\section{The extinction objection}\label{sec-extinction}

The argument so far might seem to run into a doctrine that Broome states
in the same year as the Division Principle and that his followers have
made central. When a fair lottery has been held and Michael has won the
dialysis machine, it would be absurd, Broome says, to hold another
lottery on the ground that lotteries are fair: ``once Michael has won
the first lottery he has a perfect claim to the treatment, and Maggie
has none. Consequently it is not at all fair to hold a second lottery''
(\citeproc{ref-Broome1984b}{Broome 1984b, 628--29}). Call this the
\emph{extinction thesis}: losing a fair lottery extinguishes the loser's
claim. Vong (\citeproc{ref-Vong2015}{2015}), who notes that ``there has
been little discussion of the temporal structure of how lotteries
satisfy beneficiaries' claims'' (p.~470), rejects the thesis in one form
and keeps it in another. He distinguishes a \emph{procedural} claim,
``satisfied by a lottery that gives claimants appropriate, fair chances
of benefiting,'' from a \emph{benefit} claim that ``cannot be lost due
to the results of a procedurally fair lottery'' (pp.~479--480). The
benefit claim survives losing; the procedural claim, once satisfied, is
satisfied. So when new units become available, the losers are to be
treated ``as a new case of scarce benefit distribution'' and a fresh
fair lottery run among them (p.~483). He considers, and rejects, views
on which losing a lottery weakens a claimant's claim and on which
winning strengthens it (pp.~477--479). Feldblyum Le Blevennec
(\citeproc{ref-FeldblyumLeBlevennec2023}{2023}) proposes a ``weighted
diachronic lottery'' for cases in which new claimants arrive after a
lottery has been held but before the good is handed over; but in her
scheme, as in a proposal of Broome's that Vong reports (p.~481), the
losers of the first lottery are \emph{excluded} from the second, and it
is the winner whose chances are protected.

The objection comes in two forms. On Broome's, the persistent loser in
New Mexico has no more claim than anyone else: each year her claim was
fairly entered, surrogately satisfied by an equal chance, and
extinguished. On Vong's, her benefit claim survives, but so does
everyone else's, equally, and what she is owed each year is a fresh fair
lottery among all who remain, which is the forgetting rule. On neither
is there anything for memory to track, and the existing literature on
lotteries with memory, such as it is, has memory running the other way:
it is winners whose history counts.

The reply is to distinguish three structures that the objection runs
together.

The first is a \emph{re-run}: the same claimants, the same unit, a
second lottery held because the first went the wrong way for someone.
Broome is right that this is absurd, and the reason is instructive:
there is nothing left to divide. The loser's claim to \emph{that} unit
was given its surrogate and the unit is now the winner's; another
lottery would take it back from someone with a perfect claim to it. The
extinction thesis is correct about the re-run, and I accept it.

The second is the \emph{interrupted allocation} that Vong and Feldblyum
Le Blevennec discuss: one unit, a lottery held, and then new claimants
who arrive before the unit is handed over. Here too there is one unit,
and memory of \emph{winning} is what the case calls for, since the
problem is to protect the standing of a winner whose lottery was fair
without treating late arrivals as though they had no claim. I have no
quarrel with this literature.\footnote{Vong
  (\citeproc{ref-Vong2012}{2012, 9--10}) himself appeals to repetition,
  observing that when the same two groups meet in a conflict every
  month, always saving the greater number lets the smaller group lose
  every time, and taking this to favour a lottery over majoritarianism.
  That is an argument from repetition to lotteries; mine is one from
  repetition to memory in lotteries, and they are compatible.}

The third structure is \emph{recurrence}: a new unit each period,
allocated among claimants who overlap from period to period. In this
structure the extinction thesis is true and irrelevant. Maggie's claim
to the 2026 tag was extinguished when she lost the 2026 lottery; it is
someone else's. But there is a 2027 tag, and the Division Principle says
that the 2026 and 2027 tags together are a divisible good, so that the
2027 tag goes, or is more likely to go, to someone who did not get the
2026 one. Maggie's claim to the 2027 tag is not the ghost of her claim
to the 2026 tag. It is a claim to a unit that exists, and it is stronger
than the claim of someone who received the 2026 unit for the same reason
that, when a cake is being divided, the claim of someone who has had no
slice yet is stronger than the claim of someone who has had one. Where
the currency is waiting, the same holds with waiting in place of slices:
her claim to be served now is stronger than a newcomer's because she has
borne more of the delay being divided.

So the extinction thesis and the memory thesis are about different
goods. What Broome's form of the objection needs, in order to bite, is a
stronger thesis: that losing a fair lottery for one unit extinguishes
not only the claim to that unit but any bearing the loss might have on
the allocation of later units. And that stronger thesis is what the cake
refutes. Someone who raffled each slice separately could say, of each
raffle, that the losers' claims to that slice had been fairly dealt with
and extinguished; it would still be true that the cake had not been
divided. Vong's form fails for a related reason. A procedural claim is a
claim to a fair procedure, and it is satisfied by a fresh equal lottery
only if a fresh equal lottery is the fair procedure; over recurring
units, the Division Principle says it is not. The option Vong does not
consider, between the weakening and strengthening views he rejects, is
that losing strengthens the claimant's claim to \emph{later} units; and
his objections to those views, which turn on the arrival of new
claimants for a single unit, do not touch it. His own case against
Broome is, in fact, a demand for memory of a modest kind: the time at
which one fell overboard should not affect one's cumulative chance of
the buoy, so later lotteries for one unit must be adjusted for earlier
ones (pp.~474--476). The argument of this paper extends that demand from
one unit to the sequence.

A residual worry is worth drawing out, because answering it is what
makes the surrogate character of chances do real work. If the lottery
gave Maggie surrogate satisfaction in 2026, then in 2026 her claim
\emph{was} satisfied, if only by surrogate; is extra weight in 2027 not
satisfying it twice? This is where it matters that a surrogate is a
surrogate. A chance is not a portion of the good; it is what one gives
when a portion cannot be given. If it were a portion, so that an equal
chance of a slice were a kind of equal slice, then the raffle of the
cake would be a division of the cake, and the theory could not say, as
it does, that the cake should be divided instead. The whole content of
the Division Principle is that surrogate satisfaction is not to be
preferred to real satisfaction when the real thing is available, and in
the recurrence structure it is available over the sequence. So the
surrogate Maggie received in 2026 is no reason to withhold from her the
2027 unit when it can be given.\footnote{The structure resembles what
  Thomson (\citeproc{ref-Thomson1990}{1990}) calls the residue of a
  justified infringement, with one difference that matters in
  Section~\ref{sec-reach}: nothing was infringed. The lottery was the
  fairest thing available. What remains after it is not a wrong to be
  compensated but a division still to be completed. This also answers
  the objection that memory ``double counts'' chances. What is counted,
  on the view I am defending, is not chances but units received. Chances
  are the currency in which we pay when we cannot pay in units; they are
  not themselves entered in the ledger. Wasserman's
  (\citeproc{ref-Wasserman1996}{1996}) remark that one cannot eat
  chances is, on this view, the beginning of an argument for memory
  rather than an objection to lotteries.}

\section{What memory tracks}\label{sec-proportion}

\subsection{Two structures}\label{two-structures}

Sort the institutions of Section~\ref{sec-intro} by structure before
asking what fairness requires of them. In the \emph{recurrence}
structure, a claimant who receives a unit stays in the pool and may
receive another: hunting tags, river permits, chores, ticket ballots,
the annual lotteries some funders now use for research grants. In the
\emph{bottleneck} structure, a claimant who receives a unit leaves: a
visa, a flat, a kidney, a school place. The structures are individuated
by whether winners leave, the currency of memory by Fungibility, and the
two are independent; but they run together in practice, since the
claimant who leaves on being served was waiting for one unit and cared
when it came. Mixed cases, a rota whose nights are not interchangeable,
a grant that comes too late for the laboratory, have units discounted by
delay as their currency; and a bottleneck whose claimants did not care
when they were served would be one where no one is waiting for anything,
the forgetting rule is fair, and the persistent loser has no complaint.
I take the two pure cases in turn.

The designs in use track three things. \emph{Points for losses} track
how often a claimant has lost since she last won, and raise her chances
linearly (Arizona), convexly (Nevada), or in reverse, by lowering the
chances of recent winners (the Grand Canyon). \emph{Queues} track how
long a claimant has waited and serve whoever has waited longest, chance
only breaking ties (Colorado, Wyoming, and in its principal term kidney
allocation). \emph{Caps} stop a claimant at a limit. And they differ in
how hard they press what they track: a queue makes history decisive, a
linear point system makes it a matter of degree.

\subsection{Shortfall in units}\label{shortfall-in-units}

In the recurrence structure the argument of Section~\ref{sec-division}
fixes what memory should track. Fairness requires that the sequence be
divided in proportion to claims; the measure of a claimant's treatment
is the difference between the share of the units she should have
received by now and the number she has received. Call this her
\emph{shortfall}. With equal claims and one unit per period, a claimant
present for \(r\) periods who has received \(w\) units has a shortfall
of \(r/N - w\). It can be negative; a claimant who has received more
than her share has a surplus. It is measured in units, and it is signed.

Shortfall is neither the number of losses nor time waited.
Points-for-losses systems reset to zero on a win, so that a claimant who
has won once in twenty years is treated as a newcomer though her
shortfall is large, and they count a loss in a year with two hundred
applicants exactly as a loss in a year with two. Time waited tracks
shortfall only if the numbers of units and claimants are constant. What
both get right, against the forgetting rule, is that a claimant's past
bears on her present claim; what they get wrong is the quantity.

Consider, then, the family of rules that track shortfall directly and
press it with a weight \(\kappa \geq 0\): at each period, give each
claimant a chance proportional to \(1 + \kappa \cdot \text{shortfall}\),
with any negative value set to zero. At \(\kappa = 0\) this is the
forgetting rule. As \(\kappa\) grows, a claimant's chance is
increasingly dominated by her shortfall, and in the limit the rule gives
the unit to whoever has the largest shortfall, which with equal claims
is exclusion of past winners with random tie-breaking. At
\(\kappa = 1\), a claimant's expected shortfall shrinks by a factor of
\(1 - 1/N\) each period, so it is corrected over roughly the time the
practice takes to hand out one round of shares.

The floor at zero is the mechanism behind the third claim. A claimant
who has received more than her share cannot be given a negative chance;
the most one can do is exclude her until the others catch up. This is
the extinction thesis built into the rule: past winners keep what they
won. But it introduces an asymmetry: shortfalls can be pressed upward
without limit, surpluses only down to zero, and the more heavily
\(\kappa\) presses, the more the clipped surpluses fail to offset the
amplified shortfalls. In a fixed population this is harmless. Everyone
starts level, so every claimant has the same chance of every unit
whatever \(\kappa\) is, by symmetry, and as \(\kappa\) rises the
dispersion of totals falls: with ten claimants and ten periods, the
expected number left with nothing falls from 3.5 under the forgetting
rule to 2.1 at \(\kappa = 1\), 1.3 at \(\kappa = 2\), and 0 at
\(\kappa = 10\). In a fixed population, then, the Division Principle
says: press the shortfall as hard as you can, and rotate. The asymmetry
bites only when claimants come and go.

\subsection{Newcomers}\label{newcomers}

The population is not fixed. Hunters take up the sport and abandon it;
applicant pools turn over; households form and dissolve. And once
claimants come and go, a real objection to memory appears that the
fixed-population case hides. In the year in which she first applies, a
newcomer has a claim exactly as strong as anyone else's. Under any rule
with memory, she is given a smaller chance than an incumbent who has
lost before. So the rule gives unequal chances to equal claims, which is
what the standard theory says unfairness \emph{is}. Memory, on this
objection, corrects one unfairness by committing another, and the
newcomer pays.

The objection has an answer, and the answer also fixes how heavily
memory may press. Return to the two things I drew from the standard
theory in Section~\ref{sec-broome}: division is required where it is
available, and where it is not, the surrogate is owed. Division of the
sequence is available only to a claimant who stays long enough for it to
be completed. A claimant who leaves after three periods cannot be given
three-tenths of a tag; for her, the sequence is as indivisible as the
unit. If the allocator knew who would leave, it would owe such claimants
the surrogate outright: a chance, in each period, proportional to their
claim. It does not know. So it owes the surrogate to everyone, in the
only form in which it can be given without knowing tenures in advance:
that no claimant's expected share, over whatever tenure she turns out to
have, fall short of proportionality to her claims. Call this the
\emph{surrogate constraint}. It is what the second-best status of
lotteries requires when the allocator cannot tell who will stay long
enough to be treated by division, and it is prior to the Division
Principle's demand for the same reason the surrogate is owed at all:
fairness may reduce the unfairness of the sequence's outcome only by
means that do not deny anyone the chance her claim warrants. What
fairness requires is the rule that presses shortfall as hard as the
surrogate constraint allows.

Which rules those are is a matter of calculation. Suppose claimants
leave with a fixed probability each period, independently of their
history, so that tenure is random and unknown to claimant and allocator
alike, with a mean of ten periods, and the pool is kept at ten by
replacing those who leave. The question is what a claimant's expected
total is, as a function of how long her tenure turns out to be. To ask
this is to take the claimant over her whole tenure, rather than at a
period, as the unit fairness compares: the analogue of the ``complete
lives'' view of equality that McKerlie
(\citeproc{ref-McKerlie1989}{1989}) distinguishes from views comparing
simultaneous segments of lives (see also
\citeproc{ref-Temkin1993}{Temkin 1993}, ch.~8), forced here because
claims belong to persons, who persist through the periods, not to
persons-at-periods.

Table~\ref{tbl-churn} reports the results. Under the forgetting rule,
expected share is proportional to tenure at every tenure, as it must be,
since each period gives each claimant present exactly \(1/N\); its only
defect is the dispersion of realized shortfalls, which is at its
largest. (This fixes where the argument with the ex ante theorist has to
take place: the forgetting rule's defect is entirely ex post, and
Section~\ref{sec-exante} addresses that theorist directly.) Under the
shortfall rule at \(\kappa = 1\), expected share stays within two or
three per cent of proportional at every tenure, because a newcomer's
shortfall is zero, the shortfalls of the incumbents around her average
zero, and surpluses are clipped in only about one claimant-period in
thirty; and the dispersion of realized shortfalls falls by two-fifths.
This rule satisfies the surrogate constraint, to the accuracy the
simulation can resolve. At \(\kappa = 2\) the newcomer's chance is 97
per cent of what her claim warrants, a small but real breach, and
dispersion falls by half. At \(\kappa = 10\), which is what one gets by
pressing the corrective claim at full strength against the very next
unit, the newcomer's chance is little more than half of what her claim
warrants.

\begin{longtable}[]{@{}
  >{\raggedright\arraybackslash}p{(\linewidth - 12\tabcolsep) * \real{0.3171}}
  >{\raggedright\arraybackslash}p{(\linewidth - 12\tabcolsep) * \real{0.1220}}
  >{\raggedright\arraybackslash}p{(\linewidth - 12\tabcolsep) * \real{0.1220}}
  >{\raggedright\arraybackslash}p{(\linewidth - 12\tabcolsep) * \real{0.0976}}
  >{\raggedright\arraybackslash}p{(\linewidth - 12\tabcolsep) * \real{0.0976}}
  >{\raggedright\arraybackslash}p{(\linewidth - 12\tabcolsep) * \real{0.0976}}
  >{\raggedright\arraybackslash}p{(\linewidth - 12\tabcolsep) * \real{0.1463}}@{}}
\caption{Ten claimants with equal claims, one unit per period, each
claimant leaving with probability 0.1 per period and replaced. Chances
and expected shares are multiples of the proportional value (\(1/N\) per
period); ± gives one standard error across runs (40 runs; 120 for
\(\kappa = 1\); the first column's standard errors are below 0.005).
Cells for tenure 1, 3, 10 and 30 are based on roughly 5,100, 4,200,
2,000 and 230 leavers per 40 runs. Dispersion is the standard deviation,
in units, of a claimant's shortfall at the moment she leaves. Details
and sensitivity checks in the Appendix.}\label{tbl-churn}\tabularnewline
\toprule\noalign{}
\begin{minipage}[b]{\linewidth}\raggedright
Rule
\end{minipage} & \begin{minipage}[b]{\linewidth}\raggedright
Newcomer's first-period chance
\end{minipage} & \begin{minipage}[b]{\linewidth}\raggedright
Expected share, tenure 1
\end{minipage} & \begin{minipage}[b]{\linewidth}\raggedright
Tenure 3
\end{minipage} & \begin{minipage}[b]{\linewidth}\raggedright
Tenure 10
\end{minipage} & \begin{minipage}[b]{\linewidth}\raggedright
Tenure 30
\end{minipage} & \begin{minipage}[b]{\linewidth}\raggedright
Dispersion of realized shortfall
\end{minipage} \\
\midrule\noalign{}
\endfirsthead
\toprule\noalign{}
\begin{minipage}[b]{\linewidth}\raggedright
Rule
\end{minipage} & \begin{minipage}[b]{\linewidth}\raggedright
Newcomer's first-period chance
\end{minipage} & \begin{minipage}[b]{\linewidth}\raggedright
Expected share, tenure 1
\end{minipage} & \begin{minipage}[b]{\linewidth}\raggedright
Tenure 3
\end{minipage} & \begin{minipage}[b]{\linewidth}\raggedright
Tenure 10
\end{minipage} & \begin{minipage}[b]{\linewidth}\raggedright
Tenure 30
\end{minipage} & \begin{minipage}[b]{\linewidth}\raggedright
Dispersion of realized shortfall
\end{minipage} \\
\midrule\noalign{}
\endhead
\bottomrule\noalign{}
\endlastfoot
Forgetting (\(\kappa = 0\)) & 1.00 & 0.98 ± .04 & 1.03 ± .03 & 1.02 ±
.02 & 0.98 ± .04 & 0.95 \\
Shortfall, \(\kappa = 1\) & 1.01 & 0.96 ± .02 & 0.99 ± .02 & 1.02 ± .01
& 1.02 ± .01 & 0.58 \\
Shortfall, \(\kappa = 2\) & 0.97 & 0.97 ± .03 & 1.00 ± .03 & 0.99 ± .01
& 1.02 ± .02 & 0.49 \\
Shortfall, \(\kappa = 10\) & 0.56 & 0.56 ± .03 & 0.96 ± .02 & 0.94 ± .01
& 0.97 ± .01 & 0.41 \\
Points for losses, linear (Arizona) & 0.26 & 0.25 ± .02 & 0.48 ± .02 &
0.97 ± .02 & 1.11 ± .03 & 0.70 \\
Points for losses, squared (Nevada) & 0.08 & 0.08 ± .01 & 0.23 ± .01 &
1.01 ± .01 & 1.22 ± .02 & 0.58 \\
Queue / preference points (Colorado) & 0.00 & 0.00 & 0.00 & 1.01 & 1.28
± .01 & 0.46 \\
\end{longtable}

The rules in use breach the constraint outright. A linear point system
gives a newcomer a quarter of the chance her claim warrants and a
claimant who happens to stay thirty periods 11 per cent more than her
share; squared points cut the newcomer to a twelfth; a queue gives her
nothing, and anyone who leaves within a few periods leaves with nothing.
With a larger pool the effect is more severe: with a hundred claimants
and a mean tenure of twenty periods, linear points give a newcomer 7 per
cent of her proportional chance and a queue gives anyone with less than
thirty periods' tenure essentially nothing (Appendix). These rules
reduce dispersion, some by more than the moderate shortfall rules, but
they do it by transferring expected share from claimants whose tenure
turns out short to claimants whose tenure turns out long, which is to
deny the short-stayers the surrogate they are owed.\footnote{Huesch and
  Brady's subjects, who favoured excluding the first winner, were
  designing a raffle among a fixed ten with no newcomers, where
  exclusion is what the Division Principle recommends; their verdict is
  not evidence for exclusion where the pool turns over.}

So, in the recurrence structure, the third claim is established. Memory
is required; the forgetting rule is dominated at no cost in anyone's
expected share. But memory may press only as hard as the surrogate
constraint allows; to remember lexically is to let an incumbent's
accumulated claim override a newcomer's present one, and ``weaker claims
must not simply be overridden by stronger ones.'' When the pool turns
over, fairness requires a rule that raises a claimant's chances with her
shortfall at a rate that corrects it over roughly the period in which
the practice would ordinarily deliver her share; the region around
\(\kappa = 1\) has this property, and the rules in use lie well outside
it. It should be said plainly what this achieves: with the pool turning
over, a proportionate rule roughly halves the dispersion of outcomes the
forgetting rule produces, and a residual lottery remains. That is not a
defeat. The residual is exactly the surrogate the short-stayers are
owed, and a rule that eliminated it would take from them what fairness
requires they have. Only where the pool is fixed, and everyone will be
present for the division, is the residual zero, and there the rule is
rotation.

\subsection{Bottlenecks and waiting}\label{sec-bottleneck}

In the bottleneck structure, where winners leave, nobody accumulates
more than one unit, and what varies from claimant to claimant is not
units received but time waited. Fungibility fails, because each
claimant's claim is to one unit and to having it soon. What is divisible
is the waiting.

Under the forgetting rule, waiting is itself raffled. With a pool of ten
kept full by arrivals and one unit per period, some are served at once
and about one in seventy waits forty periods or more, though everyone's
expected wait is nine. A queue serves everyone at nine. Waiting is a
bad, claims to avoid it are equal, and time, unlike a flat, is
divisible. So the Division Principle applies to the waiting, and it
says: equalize it. This is why John and Millum
(\citeproc{ref-JohnMillum2020}{2020}) are right that a queue is often
fairer than a lottery, and why Hersch and Rowe
(\citeproc{ref-HerschRowe2024}{2024}) are right to distinguish the
``bottleneck'' cases where queues are appropriate from the ``scarcity''
cases where lotteries are: in a bottleneck everyone will be served and
what is divisible is the delay; in a one-shot scarcity nothing is
divisible; between them lies the recurrence structure, where what is
divisible is the units. It is also why the queue is the right kind of
rule in the bottleneck structure though the wrong kind in the recurrence
structure. There a claimant's claim is to a share of the units, accruing
period by period, and a queue gives a short-stayer none of hers; here
her claim is to one unit and to not being kept waiting longer than
others, and, where everyone will be served, a queue gives everyone
exactly that. The newcomer at the back is not being overridden; her
claim to be served \emph{now} would be a claim to jump the queue.

The surrogate constraint carries over, and it qualifies the queue as it
qualified the shortfall rule. Where everyone will be served, the
division of the waiting is available to all, and the queue completes it.
But where claimants may leave unserved, by dying, losing eligibility, or
giving up, and the allocator cannot tell who, those to whom the division
is not available are owed the surrogate: a chance, in each period,
proportional to their claim. A pure queue denies it: everyone who leaves
unserved leaves having been given no chance at all. So the fair rule in
a bottleneck with unforeseeable attrition is a waiting-weighted lottery,
in which chances rise with time waited at a rate that keeps every
claimant's expected treatment proportional over whatever stay she turns
out to have. This is the bottleneck counterpart of the moderate
shortfall rule, and several real systems are of this kind: the Grand
Canyon's tickets-for-years and Arizona's extra numbers are
waiting-weighted lotteries, not queues. Only where exit is negligible is
the pure queue the fair rule. And where an exit is foreseeable and
imminent, the claim is urgent, and whatever rule is used must be
weighted for it, as transplant allocation is; where claims differ in
strength, waiting should likewise be weighted by claim.\footnote{The
  economics of dynamic allocation reaches related conclusions by a
  different route: Leshno (\citeproc{ref-Leshno2022}{2022}) shows that
  in an overloaded waiting list a simple randomized assignment policy
  can reduce the misallocation that strict waiting-time priority
  produces, and Arnosti and Shi (\citeproc{ref-ArnostiShi2020}{2020})
  compare lotteries and wait-lists for affordable housing. Those results
  concern efficiency; the argument here concerns what the applicants are
  owed.}

The verdict for the visa lotteries follows. They are bottlenecks with
attrition: a selected worker leaves the pool, and an unselected one may
lose her sponsor or her status. The fair rule is a lottery weighted by
years of unsuccessful registration, and by urgency and the strength of
claims. It is not the forgetting rule, and it is not the rule that took
effect in 2026, which weights registrations by the wage level of the
job, a reason of usefulness rather than a claim
(Section~\ref{sec-broome}): that reform made the lottery more responsive
to the employer's interest and left it blind to the applicant's history.
The quantity to weight by is on the record the government already keeps:
the number of years each unique beneficiary has registered and been
passed over. Hersch and Rowe (\citeproc{ref-HerschRowe2024}{2024, secs.
1, 3}) reach the opposite verdict: since demand for visas will never be
met, they classify the case as scarcity rather than bottleneck and
endorse the lottery. But their categories are built for a single
allocation, and they say themselves that in a bottleneck a periodic
lottery is unfair because a person ``who arrives into the lottery pool
early and is continuously unlucky'' may wait far longer than a lucky
newcomer (sec.~4). That is true whether or not everyone will eventually
be served; what makes it true is that the same claimants meet the same
allocator year after year, and the visa lotteries have that feature
whichever category they fall in.

\section{How far does memory reach?}\label{sec-reach}

Any view on which past allocations bear on present claims must say how
far the bearing extends. If Nevada must remember Maggie's 2026 loss in
2027, must it remember her loss in the housing lottery? Her loss in
Arizona's draw? Her mother's? A view that cannot stop somewhere
principled collapses into a general demand that luck be equalized, which
is a different and far more contested thesis. In this section I derive
the bounds on memory from the argument that generates it. There are
four.

\subsection{Memory is bound to the
practice}\label{memory-is-bound-to-the-practice}

The corrective claim of Section~\ref{sec-extinction} arises from a
specific fact: that a particular agency's allocation of a particular
kind of good left residual unfairness because its units could not be
divided. The claim is a claim on the division of \emph{those} units,
addressed to \emph{that} agency, and both restrictions do work.

The good matters because the 2027 tag can be used to divide the sequence
of tags only because it is a tag. It cannot be used to divide the
sequence of flats, because a flat is not a portion of the good the
hunter has a claim to. The Division Principle requires that a divisible
good be divided; it does not require that shortfalls in one good be made
up in another. So a loss in the housing lottery generates no claim in
the tag lottery. Goods are individuated by claims, not by names: a
Nevada sheep tag and an Arizona sheep tag are different goods, because a
claim to hunt in one state is not satisfied by a tag for the other,
whereas two lotteries run by one government for the same visa under
different labels would be one good, and the shortfall would carry across
them. An allocator cannot shed its obligation by splitting a practice in
two, because the claims that fix what the good is do not split with it.

The agency matters because the corrective claim is a claim that a
division be completed, and a division is something an agent does with
what it has. Where two independent allocators divide units of one kind
among overlapping claimants, each is bound by the division it is making
and neither by the other's: a claimant's shortfall with the first is a
fact about how the first has treated her, and the second has none of the
first's units to give. The case is narrow, since claims rarely range
indifferently over two allocators' units; but where it arises the bound
holds, because the Division Principle constrains the divider.

This is what distinguishes the thesis from luck egalitarianism, which
holds that it is bad, or unjust, for one person to be worse off than
another through no fault of her own (\citeproc{ref-Dworkin1981}{Dworkin
1981}; \citeproc{ref-Cohen1989}{Cohen 1989};
\citeproc{ref-Temkin1993}{Temkin 1993}). That thesis concerns welfare or
resources in general, and its currency is compensation: a shortfall in
one good can be made up in another. Mine concerns the fair division of a
specific good by the agency dividing it, in the good's own currency, and
says nothing about whether the persistent loser is worse off overall.
Nor is she a gambler whose losses are her own affair, in Dworkin's
sense: she did not choose to gamble; the lottery was imposed as a second
best to the division fairness required, and its results are the residue
of that second best, not option luck.

\subsection{Memory is bound to the
claim}\label{memory-is-bound-to-the-claim}

The corrective claim is a strengthened claim to the good, not a credit
the claimant holds. It lasts exactly as long as her claim to the good
lasts. Someone who applied unsuccessfully for a tag for ten years and
then gave up hunting is owed nothing; there is no tag she has a claim
to, so there is nothing her shortfall could raise her chances at. Memory
is not compensation. Compensation is what one owes when a wrong has been
done and its effects must be undone; what the lottery owes its
persistent losers is a fair division of the good it is still dividing.

Four questions about the claim's boundaries have determinate answers on
this view. \emph{Transition}: when a forgetting practice adopts memory,
do losses under the old rule count? Yes. A claimant's shortfall existed
whether or not the rule tracked it, and a practice that begins to track
it has no principled reason to start the count at zero; this is the
opposite of what a rectificatory view would say, since under the old
rule no one did anything wrong. \emph{Absence}: does a period in which a
claimant did not apply count toward her share? No.~She asserted no claim
on that period's unit, so her share to date is the sum over the periods
in which she did; Nevada's rule that points lapse after two years of not
applying is, in this respect, right. \emph{Time-varying claims}: what is
a claimant's shortfall when her claim was weak in 2026 and strong in
2028? Her share to date is the sum, over the periods in which she had a
claim, of her claim in that period as a fraction of all claims present,
and her shortfall is that sum less what she received; this is how
weighting by the strength of claims and memory combine, since both are
inputs to one quantity. \emph{Regress}: if the 2027 allocation must
correct the 2026 shortfall, must the 2028 allocation correct the 2027
allocation's failure to correct it fully? No.~Shortfall is a state, not
a stack of debts: at each period there is one quantity, the difference
between what the claimant has received and her share to date, and
correcting it is the whole of what is owed.

\subsection{Memory is bound to the
claimant}\label{memory-is-bound-to-the-claimant}

Whose history? The claimant's, and the claimant is by default the
person, because claims are duties owed to someone and a person is who
they are owed to. This is not a matter the practice may settle as it
likes: New Mexico cannot make its persistent losers vanish by declaring
that its claimants are ``this year's applicants,'' a new set each year,
since the person who applied last year and this is one person with one
claim, and it is she to whom the undertaking was made. What a practice
may do is pool the claims of several persons into one claimant, as
housing systems do when they register households and as hereditary
allotments once did with lineages; then the pooled claimant's history is
the pool's, and persists as the pool does, through changes in its
membership. I accept the implication for lineages. It follows from a
decision about whose claims are pooled, not from any licence to define
persisting claimants out of existence.

What the thesis does not say is that there is a general obligation to
correct the shortfalls of one's ancestors. Claims of that kind, where
they exist, are claims of rectification, and they arise from wrongs:
takings, exclusions, breaches. The argument of this paper begins from
the Division Principle, which governs the fair division of a good one is
dividing, and goes no further than that principle goes. A practice that
divided its good fairly, by rotation, owes nothing to anyone's
descendants; a practice that raffled it owes its present claimants a
division of what remains; and a practice that stole it owes something
else, on grounds this paper does not supply.

\subsection{Memory is bound in weight}\label{memory-is-bound-in-weight}

The fourth bound is the one Section~\ref{sec-proportion} established,
and it can be stated in a sentence. Memory may press a claimant's
shortfall only as hard as the surrogate constraint allows: never so hard
that anyone who may leave before the division is complete is denied the
chance her present claim warrants. Where everyone will stay, that is no
bound at all, and the rule is rotation or the queue; where the pool
turns over, it is the bound that separates the moderate rules from the
lexical ones.

\section{Objections}\label{sec-objections}

\subsection{Losing on purpose}\label{losing-on-purpose}

Every allocator that has considered memory has met the objection that it
rewards losing, and the objection is serious. The NBA's draft lottery
gives the worst teams the best chance at the best players, and the
predictable result was that teams lost on purpose; the league flattened
the odds in 2019, so that the three worst records now share a 14 per
cent chance at the first pick rather than the worst holding 25 per cent,
for exactly that reason.\footnote{NBA, ``NBA Draft Lottery: How it
  works'' (2026). Elster (\citeproc{ref-Elster1988}{1988}) noted the
  incentive problem with the earlier version of the draft.} If Nevada's
squared points are generous enough, someone might enter the sheep draw
for years with no intention of hunting sheep, to build points. If the
visa lottery counted years of registration, employers might register
workers they had no intention of hiring.

Two things should be said, and one conceded. First, the objection
presupposes that losing is under the claimant's control, and in the
cases that motivate this paper it is not: an applicant cannot make
herself lose; she can only apply. What she can do is apply without a
genuine claim, and that is the problem of verifying claims, which every
allocation faces. Second, where losing \emph{is} under the claimant's
control, as in a sports league, the objection shows that fairness is in
tension with incentive, which is a reason to weigh fairness against
other things, not a reason to think fairness itself requires forgetting;
the NBA did not conclude that a memoryless draft would be fairer, but
that some fairness was worth trading for competitive integrity. What
must be conceded is that memory raises the value of a manufactured
claim: under the forgetting rule a speculative entrant gains one ticket,
under a rule with memory a ticket that grows. A practice that cannot
verify claims therefore has more to lose from memory than one that can,
and the right response is to attach memory to participation that costs
something, as Nevada does by requiring a licence to be bought each year,
and as the visa lottery could do through the registration fee it already
charges each employer for each beneficiary, \$215 since 2025. This is a
constraint on implementing fairness, sometimes a serious one; it is not
evidence that fairness has no memory.

\subsection{Expression}\label{expression}

Wasserman (\citeproc{ref-Wasserman1996}{1996}), Stone
(\citeproc{ref-Stone2011}{2011}), Burgers
(\citeproc{ref-Burgers2016}{2016}), and Spiekermann
(\citeproc{ref-Spiekermann2022}{2022}) hold, in different ways, that
what makes a lottery fair is not what it distributes but what it
expresses or excludes: equal standing, impartiality, the absence of any
reason that could justify preferring one claimant to another. On such a
view a rule with memory marks people. The reply is that past shortfall
is not a reason of the kind such views exclude. Stone's lotteries are
just because they keep illegitimate reasons out of a decision no
legitimate reason can settle; but a claimant's shortfall from her share
of the very good being allocated is the paradigm of a legitimate,
comparative reason, and a procedure that excludes it excludes a good
reason along with the bad. What the forgetting rule expresses to a
claimant who has entered it four times is not equal standing but that
the agency has kept no record of her. And Spiekermann, in setting out a
case, stipulates that a procedure whose outcome is fixed across
occasions must not be reused, since ``a repeated use of the procedure is
unfair'' (\citeproc{ref-Spiekermann2022}{Spiekermann 2022, 110} n.~2);
that is a narrower point than mine, but it grants that what happened at
the first occasion can bear on the fairness of the procedure at the
second.

\subsection{Equality, not fairness}\label{equality-not-fairness}

Someone might grant everything I have said about the pattern of outcomes
and deny that fairness is what condemns it. What is bad about the
forgetting rule, on this suggestion, is that it produces inequality, or
leaves some badly off, and those are matters for a theory of equality or
of priority to the worse off (\citeproc{ref-Parfit2000}{Parfit 2000}),
not of fairness. Hold welfare fixed, and the difference appears. Suppose
the persistent loser is, all things considered, no worse off than the
repeat winner, because she has other goods the winner lacks. Her
complaint survives. It is not that she is badly off but that this
agency, dividing this good among people with equal claims, has given her
less than her share and others more. The complaint is comparative, it
concerns the good in question, and it is addressed to the allocator: the
marks of a fairness complaint in Broome's sense. A priority view would
often recommend memory too, but for reasons that would equally recommend
giving the loser something else, and it could not explain why the claim
dies with her interest in the good or why it is owed by the allocator
rather than by anyone in a position to help.

\subsection{Ex ante}\label{sec-exante}

The last objection comes from the view, originating in Diamond's
(\citeproc{ref-Diamond1967}{1967}) reply to Harsanyi
(\citeproc{ref-Harsanyi1955}{1955}) and defended most carefully by Frick
(\citeproc{ref-Frick2015}{2015}; see also
\citeproc{ref-OtsukaVoorhoeve2009}{Otsuka and Voorhoeve 2009}), that
what matters to the fairness of a procedure is the prospects it gives
people at the time of choice, not the outcomes it produces. On that view
the forgetting rule and rotation are equally fair, because at the outset
each claimant has the same prospects under both, and the ex post
dispersion I have made so much of is how equal prospects happened to
resolve. Section~\ref{sec-proportion} conceded the premise: in a fixed
population the forgetting rule's defect is entirely ex post.

Two replies. The first is that even the ex ante view must say when the
ex ante is. If it is the beginning of the practice, the view cannot
prefer the forgetting rule to rotation, since they are ex ante
identical, and the Division Principle settles the question it leaves
open. If it is the beginning of each period, so that fairness requires
equal chances at each unit \emph{whatever has gone before}, the view is
asserting separability with respect to history-blind claims, which is
what the cake refutes. There is no third setting of the clock at which
the ex ante view both distinguishes the two rules and favours the one
that forgets.

The second reply is that the ex ante view was developed for one-shot
cases, where an outcome cannot be undone and prospects are all a
procedure can distribute: a vaccine given once, a rescue made once. In
the recurrence structure the outcome of any one period is not final,
because the sequence is still being divided, and what was a prospect in
one period becomes a fact the next period's allocation can attend to.
Nothing in the motivation for ex ante evaluation carries over to a
structure in which the allocator will have further units to give and
knows who has received what.

\section{Conclusion}\label{sec-conclusion}

Fairness has memory. What the memory tracks is a claimant's shortfall
from her fair share of the good being divided, in units where later
units are as good as earlier ones and in waiting where they are not; it
is owed by the agency dividing the good, for as long as she has a claim
to it; and it may be pressed only as hard as the surrogate every
claimant is owed allows. Each clause was derived from a single principle
the standard theory already contains, that a divisible good should be
divided, together with the observation that a good allocated repeatedly
is divisible whether or not its units are.

New Mexico's draw and the Wimbledon ballot are unfair as designed, not
because they are lotteries but because they forget; so are the visa
lotteries, and the 2026 reform made them responsive to usefulness rather
than to claims. Preference-point queues of the Colorado and Wyoming kind
are the right rule in the wrong structure. A shortfall rule is not
exotic: it is what a household does when it keeps a rota, and what a
divisor method does for parliaments, extended to units that arrive one
at a time.

The consequence for theory is that fairness, on the best account we have
of it, is not separable across occasions unless claims are allowed to
carry history, and it should have been visible from the start in the
example Broome chose to illustrate the Division Principle: military
service, divided by conscripting everyone for a short time. Everyone's
short time is a rotation. The lottery is what one does when one cannot
rotate, and one can rotate more often than the lottery's defenders have
noticed.

\section*{Appendix: simulations}\label{appendix-simulations}
\addcontentsline{toc}{section}{Appendix: simulations}

All figures reported in Section~\ref{sec-division} and
Section~\ref{sec-proportion} were obtained as follows. There are \(N\)
claimants with equal claims and one indivisible unit per period. At each
period, claimant \(i\) receives the unit with probability proportional
to a weight \(\omega_i\). The rules are: forgetting, \(\omega_i = 1\);
shortfall with parameter \(\kappa\),
\(\omega_i = \max(0,\, 1 + \kappa(r_i/N - w_i))\), where \(r_i\) is the
number of periods \(i\) has been present and \(w_i\) the number of units
she has received; linear points, \(\omega_i = 1 + \ell_i\), where
\(\ell_i\) is the number of periods since \(i\) last received a unit;
squared points, \(\omega_i = 1 + \ell_i^2\); queue, \(\omega_i = 1\) for
the claimants with the largest \(\ell_i\) and \(0\) otherwise; rotation,
a uniformly random order drawn once and then cycled.

In the fixed-population runs, the same \(N = 10\) claimants are present
for \(T = 10\) or \(T = 40\) periods; each rule was run 6,000 times for
\(T = 10\) and 3,000 times for \(T = 40\). The analytic values for the
forgetting rule agree with the simulated values to two decimal places.

In the churn runs (Table~\ref{tbl-churn}), each claimant leaves at the
end of each period with probability \(q\), independently of her history,
and is replaced by a newcomer with \(r_i = w_i = \ell_i = 0\); the
horizon is 1,500 periods with the first 200 discarded, and each rule was
run 40 times at \(N = 10\), \(q = 0.1\) (120 times for \(\kappa = 1\)).
For each claimant who leaves, her tenure, total, and shortfall are
recorded. ``Expected share, tenure \(L\)'' is the mean total among
claimants whose tenure was exactly \(L\), divided by \(L/N\); standard
errors are across runs. ``Newcomer's first-period chance'' is the mean,
over newcomers, of their probability of receiving the unit in the period
they arrive, times \(N\). For the shortfall rule this probability is
\(N/(1 + \sum_j \max(0, 1 + \kappa d_j))\), summing over the incumbents'
shortfalls \(d_j\); without the floor the sum would be
\((N-1) + \kappa \sum_j d_j\), whose expectation is \(N - 1\), giving
the newcomer \(1/N\) at every \(\kappa\). The floor raises the sum by
clipping surpluses, and the clipping grows with \(\kappa\): at
\(\kappa = 1\) about 3.6 per cent of claimant-periods are clipped, at
\(\kappa = 2\) about 14 per cent.

Sensitivity. At \(N = 10\) with \(q = 0.05\) (mean tenure 20) the
newcomer's chance, as a multiple of \(1/N\), is 1.01 at \(\kappa = 1\),
0.98 at \(\kappa = 2\), 0.55 at \(\kappa = 10\), 0.21 for linear points,
0.05 for squared points and 0 for the queue; at \(q = 0.2\) (mean tenure
5) the corresponding figures are 1.01, 0.97, 0.57, 0.36, 0.16 and 0. At
\(N = 100\) with \(q = 0.05\) (six runs of 3,000 periods) they are 1.00
at \(\kappa = 1\), 0.94 at \(\kappa = 2\), 0.07 for linear points and 0
for the queue, and under the queue expected share is 0.00 at tenures 1,
3 and 10 and 0.06 at tenure 30. The dispersion of realized shortfall
falls under every rule with memory in every configuration. The
simulation code (\texttt{simulations.py}) accompanies this paper.

\protect\phantomsection\label{refs}
\begin{CSLReferences}{1}{0}
\bibitem[\citeproctext]{ref-AleksandrovWalsh2020}
Aleksandrov, Martin, and Toby Walsh. 2020. {``Online Fair Division: A
Survey.''} In \emph{Proceedings of the AAAI Conference on Artificial
Intelligence}, 34:13557--62. 9. doi:
\href{https://doi.org/10.1609/aaai.v34i09.7081}{10.1609/aaai.v34i09.7081}.

\bibitem[\citeproctext]{ref-ArnostiShi2020}
Arnosti, Nick, and Peng Shi. 2020. {``Design of Lotteries and Wait-Lists
for Affordable Housing Allocation.''} \emph{Management Science} 66 (6):
2291--2307. doi:
\href{https://doi.org/10.1287/mnsc.2019.3311}{10.1287/mnsc.2019.3311}.

\bibitem[\citeproctext]{ref-BalinskiYoung2001}
Balinski, Michel L., and H. Peyton Young. 2001. \emph{Fair
Representation: Meeting the Ideal of One Man, One Vote}. 2nd ed.
Washington, DC: Brookings Institution Press. First edition Yale
University Press, 1982.

\bibitem[\citeproctext]{ref-Broome1984}
Broome, John. 1984a. {``Selecting People Randomly.''} \emph{Ethics} 95
(1): 38--55. doi: \href{https://doi.org/10.1086/292596}{10.1086/292596}.

\bibitem[\citeproctext]{ref-Broome1984b}
---------. 1984b. {``Uncertainty and Fairness.''} \emph{The Economic
Journal} 94 (375): 624--32. doi:
\href{https://doi.org/10.2307/2232707}{10.2307/2232707}.

\bibitem[\citeproctext]{ref-Broome1990}
---------. 1991a. {``Fairness.''} \emph{Proceedings of the Aristotelian
Society} 91: 87--101. doi:
\href{https://doi.org/10.1093/aristotelian/91.1.87}{10.1093/aristotelian/91.1.87}.

\bibitem[\citeproctext]{ref-Broome1991}
---------. 1991b. \emph{Weighing Goods: Equality, Uncertainty and Time}.
Oxford: Blackwell.

\bibitem[\citeproctext]{ref-Broome1998}
---------. 1998. {``Kamm on Fairness.''} \emph{Philosophy and
Phenomenological Research} 58 (4): 955--61.

\bibitem[\citeproctext]{ref-Burgers2016}
Burgers, Jan-Willem. 2016. {``Perspectives on the Fairness of
Lotteries.''} \emph{Res Publica} 22 (2): 209--24. doi:
\href{https://doi.org/10.1007/s11158-014-9265-7}{10.1007/s11158-014-9265-7}.

\bibitem[\citeproctext]{ref-Cohen1989}
Cohen, G. A. 1989. {``On the Currency of Egalitarian Justice.''}
\emph{Ethics} 99 (4): 906--44. doi:
\href{https://doi.org/10.1086/293126}{10.1086/293126}.

\bibitem[\citeproctext]{ref-Diamond1967}
Diamond, Peter A. 1967. {``Cardinal Welfare, Individualistic Ethics, and
Interpersonal Comparison of Utility: Comment.''} \emph{Journal of
Political Economy} 75 (5): 765--66. doi:
\href{https://doi.org/10.1086/259353}{10.1086/259353}.

\bibitem[\citeproctext]{ref-Duxbury1999}
Duxbury, Neil. 1999. \emph{Random Justice: On Lotteries and Legal
Decision-Making}. Oxford: Oxford University Press.

\bibitem[\citeproctext]{ref-Dworkin1981}
Dworkin, Ronald. 1981. {``What Is Equality? Part 2: Equality of
Resources.''} \emph{Philosophy \& Public Affairs} 10 (4): 283--345.

\bibitem[\citeproctext]{ref-Elster1988}
Elster, Jon. 1988. {``Taming Chance: Randomization in Individual and
Social Decisions.''} In \emph{The Tanner Lectures on Human Values},
edited by Sterling M. McMurrin, 9:105--80. Salt Lake City: University of
Utah Press.

\bibitem[\citeproctext]{ref-Elster1989}
---------. 1989. \emph{Solomonic Judgements: Studies in the Limitations
of Rationality}. Cambridge: Cambridge University Press.

\bibitem[\citeproctext]{ref-FeldblyumLeBlevennec2023}
Feldblyum Le Blevennec, Marie Kerguelen. 2023. {``Fairness and Chance in
Diachronic Lotteries: A Response to Vong.''} \emph{Journal of Ethics and
Social Philosophy} 26 (1): 192--203. doi:
\href{https://doi.org/10.26556/jesp.v26i1.2190}{10.26556/jesp.v26i1.2190}.

\bibitem[\citeproctext]{ref-Frick2015}
Frick, Johann. 2015. {``Contractualism and Social Risk.''}
\emph{Philosophy \& Public Affairs} 43 (3): 175--223. doi:
\href{https://doi.org/10.1111/papa.12058}{10.1111/papa.12058}.

\bibitem[\citeproctext]{ref-Goodwin2005}
Goodwin, Barbara. 2005. \emph{Justice by Lottery}. 2nd ed. Exeter:
Imprint Academic.

\bibitem[\citeproctext]{ref-Harsanyi1955}
Harsanyi, John C. 1955. {``Cardinal Welfare, Individualistic Ethics, and
Interpersonal Comparisons of Utility.''} \emph{Journal of Political
Economy} 63 (4): 309--21. doi:
\href{https://doi.org/10.1086/257678}{10.1086/257678}.

\bibitem[\citeproctext]{ref-Henning2015}
Henning, Tim. 2015. {``From Choice to Chance? Saving People, Fairness,
and Lotteries.''} \emph{The Philosophical Review} 124 (2): 169--206.
doi:
\href{https://doi.org/10.1215/00318108-2842176}{10.1215/00318108-2842176}.

\bibitem[\citeproctext]{ref-HerschRowe2024}
Hersch, Gil, and Thomas Rowe. 2024. {``Lotteries, Queues, and
Bottlenecks.''} In \emph{Oxford Studies in Political Philosophy}, edited
by David Sobel and Steven Wall, 10:186--210. Oxford University Press.
doi:
\href{https://doi.org/10.1093/oso/9780198909460.003.0008}{10.1093/oso/9780198909460.003.0008}.

\bibitem[\citeproctext]{ref-HerschRowe2025}
---------. 2025. {``Allocative Fairness.''} \emph{Philosophy Compass} 20
(5): e70042. doi:
\href{https://doi.org/10.1111/phc3.70042}{10.1111/phc3.70042}.

\bibitem[\citeproctext]{ref-Hooker2005}
Hooker, Brad. 2005. {``Fairness.''} \emph{Ethical Theory and Moral
Practice} 8 (4): 329--52. doi:
\href{https://doi.org/10.1007/s10677-005-8836-2}{10.1007/s10677-005-8836-2}.

\bibitem[\citeproctext]{ref-HueschBrady2010}
Huesch, Marco D., and Richard Brady. 2010. {``Allowing Repeat
Winners.''} \emph{Judgment and Decision Making} 5 (5): 374--79. doi:
\href{https://doi.org/10.1017/S1930297500002175}{10.1017/S1930297500002175}.

\bibitem[\citeproctext]{ref-Igarashi2024}
Igarashi, Ayumi, Martin Lackner, Oliviero Nardi, and Arianna Novaro.
2024. {``Repeated Fair Allocation of Indivisible Items.''} In
\emph{Proceedings of the AAAI Conference on Artificial Intelligence},
38:9781--89. 9. doi:
\href{https://doi.org/10.1609/aaai.v38i9.28837}{10.1609/aaai.v38i9.28837}.

\bibitem[\citeproctext]{ref-JohnMillum2020}
John, Tyler M., and Joseph Millum. 2020. {``First Come, First Served?''}
\emph{Ethics} 130 (2): 179--207. doi:
\href{https://doi.org/10.1086/705763}{10.1086/705763}.

\bibitem[\citeproctext]{ref-Kamm1993}
Kamm, F. M. 1993. \emph{Morality, Mortality, Volume i: Death and Whom to
Save from It}. New York: Oxford University Press.

\bibitem[\citeproctext]{ref-KirkpatrickEastwood2015}
Kirkpatrick, James R., and Nick Eastwood. 2015. {``Broome's Theory of
Fairness and the Problem of Quantifying the Strengths of Claims.''}
\emph{Utilitas} 27 (1): 82--91. doi:
\href{https://doi.org/10.1017/S0953820814000259}{10.1017/S0953820814000259}.

\bibitem[\citeproctext]{ref-KornhauserSager1988}
Kornhauser, Lewis A., and Lawrence G. Sager. 1988. {``Just Lotteries.''}
\emph{Social Science Information} 27 (4): 483--516. doi:
\href{https://doi.org/10.1177/053901888027004001}{10.1177/053901888027004001}.

\bibitem[\citeproctext]{ref-Lazenby2014}
Lazenby, Hugh. 2014. {``Broome on Fairness and Lotteries.''}
\emph{Utilitas} 26 (4): 331--45. doi:
\href{https://doi.org/10.1017/S0953820814000107}{10.1017/S0953820814000107}.

\bibitem[\citeproctext]{ref-Leshno2022}
Leshno, Jacob D. 2022. {``Dynamic Matching in Overloaded Waiting
Lists.''} \emph{American Economic Review} 112 (12): 3876--3910. doi:
\href{https://doi.org/10.1257/aer.20201111}{10.1257/aer.20201111}.

\bibitem[\citeproctext]{ref-McKerlie1989}
McKerlie, Dennis. 1989. {``Equality and Time.''} \emph{Ethics} 99 (3):
475--91. doi: \href{https://doi.org/10.1086/293092}{10.1086/293092}.

\bibitem[\citeproctext]{ref-OtsukaVoorhoeve2009}
Otsuka, Michael, and Alex Voorhoeve. 2009. {``Why It Matters That Some
Are Worse Off Than Others: An Argument Against the Priority View.''}
\emph{Philosophy \& Public Affairs} 37 (2): 171--99. doi:
\href{https://doi.org/10.1111/j.1088-4963.2009.01154.x}{10.1111/j.1088-4963.2009.01154.x}.

\bibitem[\citeproctext]{ref-Parfit2000}
Parfit, Derek. 2000. {``Equality or Priority?''} In \emph{The Ideal of
Equality}, edited by Matthew Clayton and Andrew Williams, 81--125.
Basingstoke: Macmillan.

\bibitem[\citeproctext]{ref-Piller2017}
Piller, Christian. 2017. {``Treating Broome Fairly.''} \emph{Utilitas}
29 (2): 214--38. doi:
\href{https://doi.org/10.1017/S0953820816000303}{10.1017/S0953820816000303}.

\bibitem[\citeproctext]{ref-Saunders2008}
Saunders, Ben. 2008. {``The Equality of Lotteries.''} \emph{Philosophy}
83 (3): 359--72. doi:
\href{https://doi.org/10.1017/S0031819108000727}{10.1017/S0031819108000727}.

\bibitem[\citeproctext]{ref-Saunders2010}
---------. 2010. {``Fairness Between Competing Claims.''} \emph{Res
Publica} 16 (1): 41--55. doi:
\href{https://doi.org/10.1007/s11158-010-9118-y}{10.1007/s11158-010-9118-y}.

\bibitem[\citeproctext]{ref-Sher1980}
Sher, George. 1980. {``What Makes a Lottery Fair?''} \emph{No{û}s} 14
(2): 203--16. doi:
\href{https://doi.org/10.2307/2214861}{10.2307/2214861}.

\bibitem[\citeproctext]{ref-Spiekermann2022}
Spiekermann, Kai. 2022. {``Good Reasons for Losers: Lottery
Justification and Social Risk.''} \emph{Economics and Philosophy} 38
(1): 108--31. doi:
\href{https://doi.org/10.1017/S0266267121000043}{10.1017/S0266267121000043}.

\bibitem[\citeproctext]{ref-Stone2007}
Stone, Peter. 2007. {``Why Lotteries Are Just.''} \emph{Journal of
Political Philosophy} 15 (3): 276--95. doi:
\href{https://doi.org/10.1111/j.1467-9760.2006.00274.x}{10.1111/j.1467-9760.2006.00274.x}.

\bibitem[\citeproctext]{ref-Stone2008}
---------. 2008. {``On Fair Lotteries.''} \emph{Social Theory and
Practice} 34 (4): 573--90. doi:
\href{https://doi.org/10.5840/soctheorpract200834431}{10.5840/soctheorpract200834431}.

\bibitem[\citeproctext]{ref-Stone2011}
---------. 2011. \emph{The Luck of the Draw: The Role of Lotteries in
Decision Making}. New York: Oxford University Press. doi:
\href{https://doi.org/10.1093/acprof:oso/9780199756100.001.0001}{10.1093/acprof:oso/9780199756100.001.0001}.

\bibitem[\citeproctext]{ref-Tank2024}
Tank, Lukas, Nils Wendler, and Jan Peter Carstensen-Mainka. 2024.
{``Absolute Fairness and Weighted Lotteries.''} \emph{Utilitas} 36 (4):
352--61. doi:
\href{https://doi.org/10.1017/S0953820824000207}{10.1017/S0953820824000207}.

\bibitem[\citeproctext]{ref-Temkin1993}
Temkin, Larry S. 1993. \emph{Inequality}. New York: Oxford University
Press.

\bibitem[\citeproctext]{ref-Thomson1990}
Thomson, Judith Jarvis. 1990. \emph{The Realm of Rights}. Cambridge, MA:
Harvard University Press.

\bibitem[\citeproctext]{ref-Timmermann2004}
Timmermann, Jens. 2004. {``The Individualist Lottery: How People Count,
but Not Their Numbers.''} \emph{Analysis} 64 (2): 106--12. doi:
\href{https://doi.org/10.1093/analys/64.2.106}{10.1093/analys/64.2.106}.

\bibitem[\citeproctext]{ref-Tomlin2012}
Tomlin, Patrick. 2012. {``On Fairness and Claims.''} \emph{Utilitas} 24
(2): 200--213. doi:
\href{https://doi.org/10.1017/S0953820812000143}{10.1017/S0953820812000143}.

\bibitem[\citeproctext]{ref-Vong2012}
Vong, Gerard. 2012. {``What to Do When There Isn't Enough: The Fair
Distribution of Scarce Goods.''} PhD thesis, University of Oxford. DPhil
thesis, submitted Michaelmas Term 2012.
\url{https://ora.ox.ac.uk/objects/uuid:56aba9fe-d3ff-4107-974b-00db9bf4e42e}.

\bibitem[\citeproctext]{ref-Vong2015}
---------. 2015. {``Fairness, Benefiting by Lottery and the Chancy
Satisfaction of Moral Claims.''} \emph{Utilitas} 27 (4): 470--86. doi:
\href{https://doi.org/10.1017/S0953820815000357}{10.1017/S0953820815000357}.

\bibitem[\citeproctext]{ref-Vong2020}
---------. 2020. {``Weighing up Weighted Lotteries: Scarcity, Overlap
Cases, and Fair Inequalities of Chance.''} \emph{Ethics} 130 (3):
320--48. doi: \href{https://doi.org/10.1086/707212}{10.1086/707212}.

\bibitem[\citeproctext]{ref-Wasserman1996}
Wasserman, David. 1996. {``Let Them Eat Chances: Probability and
Distributive Justice.''} \emph{Economics and Philosophy} 12 (1): 29--49.
doi:
\href{https://doi.org/10.1017/S0266267100003709}{10.1017/S0266267100003709}.

\bibitem[\citeproctext]{ref-WinteinHeilmann2018}
Wintein, Stefan, and Conrad Heilmann. 2018. {``Dividing the Indivisible:
Apportionment and Philosophical Theories of Fairness.''} \emph{Politics,
Philosophy \& Economics} 17 (1): 51--74. doi:
\href{https://doi.org/10.1177/1470594X17715248}{10.1177/1470594X17715248}.

\end{CSLReferences}




\end{document}
